% Глобальные стили изменяет только агент 3 после согласования.
\usepackage{iftex}
\ifLuaTeX\else
  \PackageError{simplex-report}{Use LuaLaTeX to build this report}{}
\fi

\usepackage{fontspec}
\setmainfont{Times New Roman}
\newfontfamily\cyrillicfont{Times New Roman}

\usepackage[main=russian,english]{babel}
\usepackage[autostyle=true]{csquotes}
\usepackage{microtype}
\usepackage{geometry}
\geometry{
  a4paper,
  left=30mm,
  right=10mm,
  top=20mm,
  bottom=20mm
}

\usepackage{setspace}
\let\ReportClassNormalSize\normalsize
\renewcommand{\normalsize}{\ReportClassNormalSize\fontsize{14pt}{16.8pt}\selectfont}
\onehalfspacing
\setlength{\parindent}{1.25cm}
\setlength{\parskip}{0pt}

\usepackage{amsmath,amssymb}
\usepackage{unicode-math}
\setmathfont{Latin Modern Math}
% В классе article нумерация формул сквозная по умолчанию.
\renewcommand{\theequation}{\arabic{equation}}

\usepackage{graphicx}
\graphicspath{{figures/}}
\usepackage{tikz}
\usepackage{caption}
\DeclareCaptionLabelSeparator{endash}{~--~}
\captionsetup{
  labelsep=endash,
  font=normalfont,
  skip=6pt
}
\captionsetup[figure]{name=Рисунок,position=bottom}
\captionsetup[table]{name=Таблица,position=top,justification=raggedright,
  singlelinecheck=false}
\setlength{\textfloatsep}{6pt plus 2pt minus 2pt}
\setlength{\intextsep}{6pt plus 2pt minus 2pt}

\usepackage{array,booktabs,longtable,tabularx}
\usepackage{enumitem}
\usepackage{listings}

\usepackage{titlesec}
\titleformat{\section}
  {\normalfont\bfseries\centering\fontsize{14pt}{21pt}\selectfont}
  {\thesection}{1em}{}
\titleformat{\subsection}
  {\normalfont\bfseries\fontsize{14pt}{21pt}\selectfont}
  {\thesubsection}{1em}{}
\titlespacing{\section}{0pt}{0pt}{21pt}
\titlespacing{\subsection}{0pt}{21pt}{21pt}

\usepackage{fancyhdr}
\fancypagestyle{plain}{
  \fancyhf{}
  % Компенсация разницы левого и правого полей для центра листа.
  \fancyfootoffset[L]{20mm}
  \fancyfoot[C]{\thepage}
  \renewcommand{\headrulewidth}{0pt}
  \renewcommand{\footrulewidth}{0pt}
}

\usepackage[backend=biber,style=gost-numeric,sorting=none]{biblatex}
\AtBeginBibliography{\raggedright}
\usepackage{xurl}
\usepackage[hidelinks]{hyperref}
\urlstyle{same}
\usepackage[nameinlink,noabbrev]{cleveref}

\addto\captionsrussian{%
  \renewcommand{\contentsname}{Содержание}%
}

% Заголовок приложения: \ReportAppendix{А}{Наименование}
\newcommand{\ReportAppendix}[2]{%
  \clearpage
  \refstepcounter{section}%
  \section*{ПРИЛОЖЕНИЕ #1\\#2}%
  \addcontentsline{toc}{section}{ПРИЛОЖЕНИЕ #1. #2}%
}
